Another useful property of compact operators is that they map weakly convergent sequences into strongly convergent sequences

\begin{Theorem}
\label{thm:compact-operator-weak-to-strong}
Suppose $A\in\mathcal{L}(X,Y)$ is compact and that
\begin{equation}\label{eq:compact-operator-weak-convergence-hypothesis}
x_n\rightharpoonup x\text{ (weakly) in }X.
\end{equation}
Then
\begin{equation}\label{eq:compact-operator-strong-image-convergence}
A(x_n)\to A(x)\text{ (strongly) in }Y.
\end{equation}

\end{Theorem}

\begingroup
\postdisplaypenalty=10000
\renewcommand{\endProof}{\par\nobreak\hfill$\blacksquare$\par\vspace{0.1cm}}
\begin{Proof}
Our first step will be to show that
\begin{equation}\label{eq:bounded-operator-weak-image-convergence}
A(x_n)\rightharpoonup A(x)\text{ (weakly) in }Y.
\end{equation}
Let $f\in Y^*$ be given. We must show that
\begin{equation}\label{eq:bounded-operator-weak-image-functional-limit}
\lim_{n\to\infty}f(A(x_n))=f(A(x)).
\end{equation}
To do this we define $g:X\to\mathbb{R}$ by
\begin{equation}\label{eq:bounded-operator-composed-dual-functional}
g(z)=f(A(z)),\quad z\in X.
\end{equation}
Now $g$ is linear since $f$ and $A$ are both linear, and $g$ is bounded since
\begin{equation}\label{eq:bounded-operator-composed-functional-norm-estimate}
|g(z)|=|f(A(z))|\leq\|f\|\|A(z)\|\leq\|f\|\|A\|\|z\|.
\end{equation}
Thus, $g\in X^*$, and since $x_n\rightharpoonup x$ in $X$ we have\[\begin{aligned}
\lim_{n\to\infty}f(A(x_n))&=\lim_{n\to\infty}g(x_n)\\
&=g(x)\\
&=f(A(x)).
\end{aligned}\]Since $f$ was arbitrary, $A(x_n)\rightharpoonup A(x)$.

Now suppose that $A(x_n)$ does not converge strongly to $A(x)$ in $Y$. Then there exists an $\varepsilon>0$ and a subsequence $A(x_{n_k})$ such that
\begin{equation}\label{eq:compact-operator-strong-convergence-contradiction-sequence}
\|A(x_{n_k})-A(x)\|\geq\varepsilon.
\end{equation}
Now, since $x_n$ converges weakly to $x$, so does $x_{n_k}$. Since $x_{n_k}$ is weakly convergent, it is bounded. Thus, since $A$ is compact, $A(x_{n_k})$ has a strongly convergent subsequence. However, since strong convergence implies weak convergence, and since weak limits are unique, this subsequence must converge to $A(x)$. However, this contradicts \eqref{eq:compact-operator-strong-convergence-contradiction-sequence} and completes the proof.
\end{Proof}
\endgroup

We can combine this result with Theorems~\ref{thm:sobolev-continuous-target-compact-embedding} and \ref{thm:compact-successive-sobolev-orders} to get the following corollaries.

\begin{Corollary}
\label{cor:compact-sobolev-weak-to-uniform}
Suppose that $k>m/2$ and $\Omega\subset\mathbb{R}^m$ is bounded and has the $k$-extension property. Let
\begin{equation}\label{eq:compact-sobolev-weak-to-uniform-hypothesis}
u_n\rightharpoonup u\text{ (weakly) in }H^k(\Omega).
\end{equation}
Then
\begin{equation}\label{eq:compact-sobolev-weak-to-uniform-conclusion}
u_n\to u\text{ uniformly on }\overline\Omega.
\end{equation}

\end{Corollary}

\begin{Corollary}
\label{cor:compact-sobolev-weak-to-lower-strong}
Suppose that $k$ is a non-negative integer and $\Omega\subset\mathbb{R}^m$ is bounded and has the $k+1$-extension property. Let
\begin{equation}\label{eq:compact-sobolev-successive-order-weak-hypothesis}
u_n\rightharpoonup u\text{ (weakly) in }H^{k+1}(\Omega).
\end{equation}
Then
\begin{equation}\label{eq:compact-sobolev-successive-order-strong-conclusion}
u_n\to u\text{ (strongly) in }H^k(\Omega).
\end{equation}

\end{Corollary}
